\ChapterPage{Página 2 \textbullet{} La Gaceta}{Una corrección al margen}
\begin{multicols}{2}
\section*{Del editor nocturno}
La Crónica desea rectificar la noticia de ayer. La campana del puerto no
sonó a las tres y media. Se oyó a las cuatro menos tres minutos, aunque los
tres testigos insisten en que el reloj de la plaza se detuvo a las 3:33.

La diferencia no tiene utilidad práctica. Aun así, el editor ha pasado la
mañana comparando relojes y ninguno coincide con el de la ciudad. Sus agujas
se mueven correctamente cuando se miran desde la calle. A través del cristal,
parecen señalar el suelo.

\textit{Lorem tenebris}: la oscuridad no es la ausencia de una lámpara, sino
el intervalo durante el cual la lámpara recuerda su primera llama. Esta
frase apareció en el margen de un libro de mareas fechado en 1792. Nadie en
el archivo reconoce la letra.

\section*{Objetos perdidos}
Se ha encontrado en el muelle oriental una llave de latón con el número 9.
No abre ninguna puerta de la aduana, aunque la cerradura de la entrada al
sótano ha sido pulida recientemente. Quien la reclame deberá describir los
dientes sin mirarlos directamente.

\columnbreak
\section*{Declaración del faro}
El farero Wren informa que anoche la lámpara dio una vuelta de más. El
mecanismo ha sido revisado y funciona correctamente. Wren ha pedido que las
persianas permanezcan cerradas durante el día; la junta se ha negado.

Al amanecer se veía una segunda sombra sobre el agua, bajo el faro. No
correspondía a ningún objeto situado por encima de las olas. La sombra giró
en sentido contrario a las agujas del reloj mientras el sol salía por el
este.

\section*{Anuncio}
\textsc{Morrow e Hijos} ofrece relojes, barómetros y cronómetros marinos de
confianza. Cada instrumento se prueba con tres patrones independientes y
una luna honrada. Ninguna compra incluye garantía contra mareas irregulares.

\textbf{Nota del editor:} el impresor pide que la última columna quede
vacía. Se deniega la petición. Una columna vacía en el periódico invita a
las cartas al director, y ya hemos recibido una del mar.
\end{multicols}
